
\usepackage[most]{tcolorbox}
\usepackage{xparse}

%\definecolor{teal}{HTML}{ba5536}
%\definecolor{terra}{HTML}{363237}
\definecolor{teal}{HTML}{34675c}
\definecolor{terra}{HTML}{324851}

% plain box:
\newtcolorbox{justabox}{
  enhanced,
  colback=terra!2!white,
  colframe=terra!80!white,
  boxrule=0.3mm,
  drop small lifted shadow,
  beforeafter skip=\baselineskip,
}

% Definition environment with counter:
\newtcbtheorem[number within=chapter]{definition}{Definition}{%
  enhanced,
  breakable,
  enhanced jigsaw,
  colback=terra!2!white,
  colframe=terra!80!white,
  colbacktitle=terra!2!white,
  boxrule=0.3mm,
  drop small lifted shadow,
  titlerule=-0.5mm,
  toptitle=2mm,
  beforeafter skip=\baselineskip,
  fonttitle=\upshape\bfseries\sffamily,
  coltitle=terra!90!black,
  fontupper=\normalfont,
  }{def}

% Observation environment without title bar:
\newtcbtheorem[number within=chapter]{obs}{Observation}{%
  theorem style=plain,
  enhanced,
  colback=terra!2!white,
  colframe=terra!80!white,
  boxrule=0.3mm,
  beforeafter skip=\baselineskip,
  fonttitle=\upshape\bfseries\sffamily,
  coltitle=terra!90!black,
  fontupper=\normalfont,
  drop small lifted shadow,
  }{obs}
\newenvironment{observation}[1]{\begin{obs}{}{#1}}{\end{obs}}

% Theorem environment without 'Theorem' label:
\newenvironment{theorem}[2]{%
  \begin{tcolorbox}[%
    enhanced,
    title={#1},
    colframe=terra!80,
    colback=terra!2!white,
    colframe=terra!80!white,
    colbacktitle=terra!2!white,
    coltitle=terra!90!black,
    boxrule=0.3mm,
    titlerule=-0.5mm,
    toptitle=2mm,
    fonttitle=\upshape\bfseries\sffamily,
    drop small lifted shadow,
    beforeafter skip=\baselineskip,
    ]\label{thm:#2}%
  }%
  {\end{tcolorbox}}

% Proof:
\renewtcolorbox{proof}{
  blanker,enforce breakable,left=3mm,
  before skip=10pt,after skip=10pt,
  parbox=false, % for paragraph indentation
  borderline west={0.2mm}{0pt}{terra!60!white}
}

% semantics:

\usepackage{wrapfig}
\let\description=\eqlist
\let\enddescription=\endeqlist
\let\eqlistlabel\descriptionlabel
% Formal semantics clauses. Usage:
%   \begin{semantics}
%     \clause{(a)}{M,w \models P}{$w$ is in $V(P)$.}
%     ...
%   \end{semantics}
% Each \clause takes an explicit label, a left-hand side (given without math
% delimiters), and a condition (the "iff" is supplied automatically). Long
% conditions wrap by themselves -- no manual line breaks. In the PDF the label
% and lhs columns are auto-sized so every "iff" aligns and continuation lines
% hang under the condition; in HTML the same layout is produced by a CSS grid
% (styled under .semantics in logic2.css).
\makeatletter
\newlength{\sem@labwidth}
\newlength{\sem@lhswidth}
\newlength{\sem@iffwidth}
\newlength{\sem@clausesep}
\setlength{\sem@clausesep}{0.4ex}% vertical gap between clauses
\newif\ifsem@first
\newcommand{\clause}[3]{}% placeholder, redefined inside the environment
\NewDocumentEnvironment{semantics}{+b}{%
  \settowidth{\sem@iffwidth}{iff~}%
  \ifdefined\HCode
    % ---- HTML: emit a CSS-grid structure ----
    \renewcommand{\clause}[3]{%
      \HCode{<span class="sem-lab">}##1\HCode{</span>}%
      \HCode{<span class="sem-lhs">}$##2$\HCode{</span>}%
      \HCode{<span class="sem-iff">iff</span>}%
      \HCode{<span class="sem-cond">}##3\HCode{</span>}%
    }%
    \par\EndP\HCode{<div class="semantics">}\IgnorePar
    #1%
    \HCode{</div>}\ShowPar\par
  \else
    % ---- PDF ----
    % No leading/trailing skip of our own: the surrounding source already puts a
    % \medskip/\bigskip (or a box edge) around the block. We only enter vertical
    % mode here.
    \par
    % pass 1: measure the widest label and the widest left-hand side
    \setlength{\sem@labwidth}{0pt}%
    \setlength{\sem@lhswidth}{0pt}%
    \renewcommand{\clause}[3]{%
      \settowidth{\dimen@}{##1}%
      \ifdim\dimen@>\sem@labwidth \global\sem@labwidth=\dimen@ \fi
      \settowidth{\dimen@}{\ensuremath{##2}}%
      \ifdim\dimen@>\sem@lhswidth \global\sem@lhswidth=\dimen@ \fi
    }%
    \setbox\z@=\vbox{#1}%
    \addtolength{\sem@labwidth}{0.6em}% gap between label and lhs
    \addtolength{\sem@lhswidth}{1em}%  gap between lhs and "iff"
    % pass 2: typeset. Each clause is a real paragraph; the first line holds the
    % label box, the auto-sized lhs box (so every "iff" aligns), "iff", and the
    % condition; \hangindent hangs continuation lines under the condition. A
    % \sem@clausesep skip separates clauses (but not before the first one).
    \setlength{\parindent}{0pt}%
    \setlength{\parskip}{0pt}%
    \sem@firsttrue
    \renewcommand{\clause}[3]{%
      \par
      \ifsem@first \sem@firstfalse \else \vskip\sem@clausesep \fi
      \noindent
      \hangindent\dimexpr\sem@labwidth+\sem@lhswidth+\sem@iffwidth\relax
      \hangafter\@ne
      \makebox[\sem@labwidth][l]{##1}%
      \makebox[\sem@lhswidth][l]{\ensuremath{##2}}iff~##3%
      \par
    }%
    #1%
  \fi
}{}
\makeatother


% exercises and solutions:

\ifcompilesolutions

\usepackage[clear-aux]{xsim}
\xsimsetup{
  exercise/template=tcb,
  exercise/within=chapter,
  exercise/the-counter = \thechapter.\arabic{exercise},
  solution/template=tcbsol,
  solution/name=Exercise,
  print-solutions/headings=false,
  split-aux-lists=|--INTERNAL--||
}

\DeclareExerciseEnvironmentTemplate{tcb} {%
  \tcolorbox[
  breakable,
  enhanced jigsaw,
  fonttitle=\sffamily\upshape\bfseries,
  coltitle=teal!10!white,
  colback=teal!02!white,
  colbacktitle=teal,
  colframe=teal!70!black,
  boxrule=0.3mm,
  attach boxed title to top left={yshift=-1mm, xshift=1mm},
  beforeafter skip=.9\baselineskip,
  title = Exercise~\GetExerciseProperty{counter}%
        \GetExercisePropertyT{subtitle}{ \textit{\PropertyValue}}%
        \IfInsideSolutionF{%
            \GetExercisePropertyT{points}{ % notice the space
                (%
                \PropertyValue
                \IfExerciseGoalSingularTF{points}
                    {\XSIMtranslate{point}}
                    {\XSIMtranslate{points}}%
                )%
            }%
        }%
    ]%
    % \IfInsideSolutionT{%
    %   \tcbsubtitle{Exercise}
    %   \input{\jobname-\ExerciseType-\ExerciseID-exercise-body.tex}
    %   \tcbsubtitle{Solution}
    % }
}{\endtcolorbox}

% Solution headings:
\DeclareExerciseEnvironmentTemplate{tcbsol}{%
  \tcbox[
  fontupper=\sffamily\upshape\bfseries,
  coltext=teal!10!white,
  colback=teal!80!white,
  colframe=teal!50!white,
  boxrule=0.3mm,
  boxsep=-0.5mm,
  before skip=\baselineskip,
  after skip=0.5\baselineskip,
  ]{Exercise~\GetExerciseProperty{counter}}
  \noindent
}

\else

\newtcolorbox[auto counter,number within=chapter]{exercise}{
  breakable,
  enhanced jigsaw,
  fonttitle=\sffamily\upshape\bfseries,
  coltitle=teal!10!white,
  colback=teal!02!white,
  colbacktitle=teal,
  colframe=teal!70!black,
  boxrule=0.3mm,
  attach boxed title to top left={yshift=-1mm, xshift=1mm},
  beforeafter skip=.9\baselineskip,
  title=Exercise~\thetcbcounter,
}
\usepackage{comment}
\excludecomment{solution}
\newcommand{\printsolutions}{}

\fi

% \DeclareExerciseType{exercise}{
%     exercise-env=ltxexercise,
%     solution-env=ltxsolution,
%     exercise-name=\XSIMtranslate{exercise},
%     solution-name=\XSIMtranslate{solution},
%     exercise-template=tcb,
%     solution-template=tcb,
%     within=chapter,
% }

% ========== custom list format in exercises ===========

\usepackage[shortlabels]{enumitem} % begin{enumerate}[(a)] and itemsep
\newenvironment{exlist}{\vspace{-2mm}\begin{enumerate}[(a)]\itemsep-1.5mm}{\end{enumerate}\vspace{-5mm}}
\newenvironment{sollist}{\vspace{-5mm}\begin{enumerate}[(a)]\itemsep-1.5mm}{\end{enumerate}\vspace{-3mm}}
\newcommand{\beginwithlist}{\vspace{1mm}}
\setlist[itemize]{leftmargin=*}
\setlist[enumerate]{leftmargin=*}
